\documentclass[11pt,a4paper]{article}
\usepackage[utf8]{inputenc}
\usepackage[T1]{fontenc}
\usepackage[french]{babel}
\usepackage[left=2cm,right=2cm,top=2cm,bottom=2cm]{geometry}
\usepackage[dvipsnames]{xcolor}
\usepackage{fouriernc}
\usepackage{amsfonts}
\usepackage{amssymb}
\usepackage{amsmath}
\usepackage{enumitem}
\usepackage{fancyhdr}
\usepackage{awesomebox}
\usepackage{bclogo}
\usepackage{chemmacros}
\chemsetup{modules=all}

\frenchbsetup{StandardLists=true}

\pagestyle{fancy}
\renewcommand\headrulewidth{1pt}
\fancyhead[L]{Lycée Denis Diderot}
\fancyhead[R]{Classe de 1 Bac Pro}
\fancyfoot[C]{}
\fancyfoot[R]{Année 2025-2026}
\fancyfoot[L]{Alexis RUIZ}

\newcommand{\cadreblanc}[1]{%
  \medskip\framebox[\linewidth][c]{%
  \rule{0mm}{#1mm}}\par
  \medskip}

\newcounter{numexos}
\setcounter{numexos}{0}

\newcommand{\bex}[2]{%
  \stepcounter{numexos}%
  \begin{bclogo}[couleur=white, logo=, barre=none, arrondi=0.1]%
    {\normalsize\textbf{Exercice~\thenumexos~: #1}}%
  #2
  \end{bclogo}%
}

\begin{document}

\pagenumbering{gobble}

\begin{center}
  \LARGE{Exercices}\\[3mm]
\end{center}

\hrule\
\\[0.5cm]

\bex{masse molaire}{
  À l'aide de la classification périodique, calculer les masses molaires suivantes :
  \begin{itemize}[label=\textbullet, font=\LARGE\color{black}]
    \vspace{3mm}
    \item M(\chemform{H_2O}) : \dotfill
    \vspace{3mm}
    \item M(\chemform{CO_2}) : \dotfill
    \vspace{3mm}
    \item M(\chemform{C_2H_6O}) : \dotfill
    \vspace{3mm}
    \item M(\ch{Na\pch[]}) : \dotfill
    \vspace{3mm}
    \item M(\ch{Cl\mch[]}) : \dotfill
    \vspace{3mm}
    \item M(\ch{OH\mch[]}) : \dotfill
    \vspace{3mm}
    \item M(\ch{H_3O\pch[]}) : \dotfill \\[3mm]
  \end{itemize}
}

\vfill 

\bex{préparation d'une solution}{
  On souhaite préparer 0,1~L de solution aqueuse en ions \ch{Fe\pch[3]} de concentration égale à
  0,04~mol/L à partir de cristaux de chlorure de fer~III de formule \ch{FeCl_3} :
  \[
    \ch{FeCl_3} \longrightarrow \ch{Fe^{3+}} + 3\,\ch{Cl^{-}}
  \]

  \vspace{5mm}

  \begin{enumerate}
    \item Calculer $n$, la quantité de matière en ions \ch{Fe\pch[3]} contenue dans la solution.

    \vspace{3mm}
    \dotfill

    \item Retrouver la quantité de matière de chlorure de fer, $n$, en mol, nécessaire pour
    réaliser la solution :

    \vspace{3mm}
    \dotfill

    \item Calculer $M$(\ch{FeCl_3}), en g/mol, la masse molaire du chlorure de fer~III.

    \vspace{3mm}
    \dotfill

    \item Calculer $m$, en g, la masse nécessaire pour réaliser la solution.

    \vspace{3mm}
    \dotfill \\[3mm]
  \end{enumerate}
}

\vfill
\end{document}